\section{Clustering and distances}
\label{sec:clustering}

The clustering runs average linkage under each distance and $k$-means on standardized features (the original F0 features as they are), with as many clusters as classes. Agreement between clusters and classes is near zero for every feature set, method, and distance: ARI from \N{c:ari_labels_min} to \N{c:ari_labels_max}. On BCI IV 2a, $k$-means clusters partly follow the subjects, with an ARI against subject identity up to \N{c:ari_subj_bci_max}, and against the classes at most about \N{c:ari_bci_kmeans_cls_max}. On EEG MMI and EEGMAT, two clusters can show at most a coarse grouping of 106 or 36 subjects, so the ARI against subject identity says little either way. Some configurations of the original features follow EMD row counts, such as $k$-means on the entropy features (Section~\ref{sec:baseline}). Average-linkage clustering puts at least 95\% of windows in one cluster in \N{c:avg_link_95} of \N{c:avg_link_total} configurations, over the four data set variants, BCI late included, an imbalance consistent with outliers. Their silhouettes run from \N{c:sil95_min} to \N{c:sil95_max}, and a high one there means that a small group of windows was split off, not that the classes separate. The highest, for EEG MMI Hilbert band power with the channels averaged under the Itakura--Saito distance, splits off \N{c:silmax_small} of \N{c:silmax_n} windows. In the original pipeline, the $k$-means step never used the distance, which only scored the silhouette, so the three distances' silhouettes describe one partition measured three ways and say nothing about which distance clusters better. On the original features, the zero floor of the padded rows likely shapes the Itakura--Saito distances; this was not measured.
